\documentclass[a4paper,12pt]{article}

% ------------------------------------------------
% PACOTES
% ------------------------------------------------
\usepackage[utf8]{inputenc}
\usepackage[T1]{fontenc}
\usepackage[brazil]{babel}
\usepackage{geometry}
\usepackage{array}
\usepackage{xcolor}
\usepackage{graphicx}
\usepackage{lastpage}
\usepackage{fancyhdr}
\usepackage{colortbl}

\geometry{margin=2cm}

% ------------------------------------------------
% CORES
% ------------------------------------------------
\definecolor{preto}{RGB}{0,0,0}
\definecolor{cinza}{RGB}{220,220,220}

% ------------------------------------------------
% DADOS DA PROVA (ALTERAR AQUI)
% ------------------------------------------------
\newcommand{\avaliacao}{AVALIAÇÃO 3}
\newcommand{\curso}{Engenharia de Software}
\newcommand{\disciplina}{Estrutura de Dados II}
\newcommand{\professor}{Vinicius Felix de Paula}
\newcommand{\turma}{B}
\newcommand{\campus}{Rio Verde}
\newcommand{\codigo}{ESW416}
\newcommand{\valor}{7,0}

% ------------------------------------------------
% CABEÇALHO
% ------------------------------------------------
\pagestyle{fancy}
\fancyhf{}
\rhead{Página \thepage\ de \pageref{LastPage}}
\renewcommand{\headrulewidth}{0pt}

% ------------------------------------------------
% INÍCIO DO DOCUMENTO
% ------------------------------------------------
\begin{document}

% ==================================================
% LOGO + TÍTULO
% ==================================================
\noindent
\begin{minipage}{0.25\textwidth}
    \includegraphics[width=3cm]{logo_unirv.png}
\end{minipage}
\begin{minipage}{0.55\textwidth}
    \centering
    \colorbox{black}{
        \parbox[c][1cm][c]{\linewidth}{
            \centering
            \color{white}\Large\bfseries \avaliacao
        }
    }
\end{minipage}
\begin{minipage}{0.18\textwidth}
    \centering
    \fbox{
        \parbox[c][1cm][c]{3cm}{
            \centering \textbf{NOTA}
        }
    }
\end{minipage}

\vspace{0.4cm}

% ==================================================
% TABELA INSTITUCIONAL
% ==================================================
\renewcommand{\arraystretch}{1.3}

\begin{tabular}{|p{10cm}|p{5cm}|}
\hline
\textbf{Curso} & \textbf{Campus} \\
\curso & \campus \\
\hline
\multicolumn{2}{|l|}{\textbf{Disciplina}} \\
\multicolumn{2}{|l|}{\disciplina} \\
\hline
\textbf{Nome do(a) acadêmico(a)} & \textbf{Valor} \\
& \valor \\
\hline
\textbf{Professor} & \textbf{Turma} \\
\professor & \turma \\
\hline
\textbf{Data da Avaliação} & \textbf{Código} \\
& \codigo \\
\hline
\end{tabular}

\vspace{0.3cm}

% ==================================================
% AVISO
% ==================================================
\noindent
\colorbox{black}{
\parbox{\textwidth}{
\centering
\color{white}\textbf{
ATENÇÃO: Somente serão passíveis de REVISÃO avaliações resolvidas à TINTA.
}
}
}

\vspace{0.4cm}

% ==================================================
% ORIENTAÇÕES
% ==================================================
\textbf{ORIENTAÇÕES PARA A RESOLUÇÃO}
\begin{itemize}
\item A avaliação é individual, sem consulta;
\item A interpretação de texto faz parte da avaliação;
\end{itemize}

\vspace{0.6cm}

% ==================================================
% QUESTÕES
% ==================================================
\begin{center}
\Large\textbf{QUESTÕES}
\end{center}

\vspace{1cm}

\textbf{1)} Explique o conceito de estrutura de dados.

\vspace{3cm}

\textbf{2)} Diferencie pilha e fila.

\vspace{3cm}

\end{document}  